\section{When a component reaches its next accounting event}
\label{sec:corporate-residual-transitions}

A recognised balance also promises a subsequent method. A refund liability
must eventually be settled; an impaired asset may recover; a modified lease
must continue through payments and depreciation. The following methods specify
those transitions under stated facts. They use the retained historical
IFRS text, with current authority and actual-company measurement kept as
separate admission questions \citep{eu20231803}.

\subsection{Actual returns can differ from the estimate}

Section~\ref{sec:company-composition} follows two expected and actual returns.
More generally, let \(L\) be the opening refund liability, \(A\) the recovery
asset, \(P\) the actual refund payment, \(K\) the actual selling-handling cost
and \(V\) the returned inventory admitted after deterioration and its NRV
ceiling. When all rights have closed, the debit-positive entries are
\begin{align}
 \Delta\mathrm{cash}&=-P-K,&
 \Delta\mathrm{inventory}&=V,\\
 \Delta\mathrm{refund liability}&=L,&
 \Delta\mathrm{recovery asset}&=-A,\\
 \Delta\mathrm{revenue}&=P-L,&
 \Delta\mathrm{cost of sales}&=A-V,\\
 \Delta\mathrm{selling expense}&=K.
\end{align}
The debit-positive revenue amount is negative when releasing an excessive
refund estimate increases revenue. These entries sum to zero. The separation
of selling cost follows the explicit handling classification and IAS~2.16;
IFRS~15.B25 governs the recovery asset. If handling instead creates eligible
inventory cost, that needs its own admitted classification.

With \(L=40\) and \(A=16\), one actual return pays refund 20 and handling 4,
and recovers inventory 10. Revenue increases 20, cost of sales is 6 and selling
expense 4. Three actual returns instead pay refund 60 and handling 12, recover
inventory 30 and reduce revenue by 20; cost of sales is credited 14. If two
returns recover inventory worth only 14, the inventory ceiling raises cost of
sales to 2. Unknown handling, inventory NRV or an open return window prevents
final settlement rather than supplying a zero.

\subsection{A reversal ceiling does not supply an allocation weight}

IAS~36.122--123 allocates a cash-generating unit's reversal to assets other than
goodwill, proportionately to their carrying amounts, subject to the individual
ceilings. Suppose assets A and B carry 40 and 60 and goodwill is 10. The unit's
recoverable amount rises to 170, allowing a total increase of 60. A's individual
ceiling is 45 and B's is 150. The first proportionate allocation would assign
24 to A and 36 to B. A can accept only 5; reallocating the remaining 19 to B
gives increases of 5 and 55. Goodwill remains unchanged.

Writing \(w_i\) for the original carrying amounts and \(c_i\) for the remaining
headroom, the method distributes a remaining reversal \(R\) among eligible
assets by \(Rw_i/\sum_j w_j\), caps each amount at \(c_i\), then redistributes
any remainder among assets still below their ceilings. This terminates when the
reversal is exhausted or no headroom remains. It requires a positive sum of
eligible weights whenever a positive amount remains.

If all eligible carrying amounts are zero, the denominator is zero. Neither
equal weighting nor weighting by headroom follows from the inspected source.
The specification therefore retains an unresolved allocation with unchanged
books. A different rule applies to an individual asset: a zero carrying amount,
no-impairment carrying value of 80 and recoverable amount of 120 permits an
individual reversal to 80 under IAS~36.117--119. That individual calculation
does not resolve the allocation among several zero-carrying assets.

\subsection{A lease reduction must identify the rights surrendered}

A single physical percentage is insufficient when leased rights have different
values. Consider right A with right-of-use asset 90 and liability 100, and right
B with asset 30 and liability 50. The supplied modification terminates B,
releasing asset 30 and liability 50 for a gain of 20. It leaves A with one
payment of 121 in a year at a revised 10\% rate. A's revised liability is
\(121/1.1=110\); the increase of 10 adjusts its asset from 90 to 100.
These are the scope-reduction and remeasurement steps in IFRS~16.45--46.

A further same-day reduction surrenders an admitted portion of A carrying
asset 10 and liability 11. The remaining liability before remeasurement is
99. Revised consideration of 99 due in one year has present value 90, so the
remeasurement reduces both liability and asset by 9. The final asset is
\(100-10-9=81\), liability 90 and additional gain 1. A year later, payment 99
contains interest 9 and principal 90; the supplied systematic depreciation
amount of 81 exhausts the asset.

The event must therefore carry the actual surrendered asset and liability for
each right, the revised payment schedule and the rate. The finite schedule
uses annual-end payments and a checked one-year interval. Missing attribution,
an unadmitted depreciation basis, a different calendar or a negative remaining
asset requires a distinct finding or method. It cannot be repaired by applying
the same physical fraction to every balance.

\subsection{Contract costs survive more than one service period}

IFRS~15.91--104 first distinguishes eligible contract costs from amounts governed
by another standard. Acquisition costs must be incremental and recoverable.
Fulfilment costs outside other standards must relate directly to the contract,
create or enhance resources for future performance and be recoverable. Waste,
past service and the specified administrative costs are expensed. The event
records these findings before calculating the posting.

For example, expense waste 5 and capitalise eligible fulfilment cost 100 for
five equal service units. Delivery of one unit amortises 20, leaving 80. After
other applicable asset impairments have been recognised, credit-adjusted
remaining consideration 70 less direct future cost 20 supports a contract
asset of 50: recognise impairment 30. Delivery of another unit amortises
\(50/4=12.5\), leaving 37.5. Without the impairment, the unamortised balance
would now be 60. If circumstances improve enough to support 60, reversal is
limited to \(60-37.5=22.5\); the final three units amortise 60.

That ceiling needs a separate no-impairment balance, not merely the present
asset. Subsequent CGU assessment under IFRS~15.103 remains required after the
component test. A successful component calculation does not establish that
the wider CGU has no impairment.

The same history discipline applies to repeated revenue modifications. Under
IFRS~15.21(a), eight remaining distinct units with consideration 80, followed
by four additional units for 60 that do not form a separate contract, leave
12 units and consideration 140. Delivery of three recognises 35. Adding three
more units for 30 then leaves 12 units and consideration 135, recognised on
their delivery. Earlier revenue is not reallocated.

For a single integrated obligation, IFRS~15.21(b) instead requires the
cumulative effect. A contract priced at 200 and 25\% complete has revenue 50.
Changing price to 260 and admitted progress to 20\% changes cumulative revenue
to 52, a catch-up of 2. Later 50\% progress yields 130. A further modification
to price 300 and progress 40\% reduces cumulative revenue to 120, a catch-up
of negative 10; completion ultimately brings it to 300. Arbitrary mixtures of
distinct and integrated promises require their own allocation, not a selection
of whichever result is convenient.

\subsection{A tax benefit retains the origin of the underlying item}

Acquisition-related recognition depends on information as well as dates. Under
IAS~12.68 and IFRS~3.45--46, a newly recognised acquired deferred tax benefit
from acquisition-date information within the open measurement period reduces
goodwill, with any excess recognised in profit. A benefit of 20 against
goodwill 5 therefore reduces goodwill by 5 and gives tax income 15. The
measurement period cannot exceed one year; elapsed time alone does not establish
that the new information concerns acquisition-date facts.

Other benefits retain their independently identified origin. In particular,
IAS~12.67 excludes an acquirer's own deferred tax benefit from acquisition
accounting; it does not place every such benefit in profit. IAS~12.61A--63
still routes tax associated with OCI or direct-equity items to that origin.
An acquirer-owned benefit of 20 with an OCI origin therefore adds DTA 20 and
OCI 20, leaving goodwill and profit unchanged. Missing origin remains unresolved.

Mixed tax schedules retain one lot per entity, authority, temporary difference,
recovery-period rate, recognition finding and origin. For a signed difference
\(D_i\), positive for deductible and negative for taxable, admitted rate
\(r_i\) and supported deductible recognition fraction \(f_i\), the signed
deferred amount is \(D_i r_i f_i\); taxable lots use \(f_i=1\). Changes go
to the traced profit, OCI or equity account. A supplied schedule can thus move
gross DTA from 30 to 19 and DTL from 20 to 30 while allocating expense 5 to
profit, 10 to OCI and 6 to direct equity.

Gross asset and liability amounts stay available even when presentation is
net. The finite offset route requires the same taxable entity and authority
and an enforceable right to set off current tax under IAS~12.74. It does not
implement the separate cross-entity simultaneous-settlement route. Neither a
blended rate nor an assumed right of setoff can substitute for those findings.
